\section{Related Work\label{sec:work}}

Architecture recovery and code-to-architecture mapping both infer a system's modular structure from implementation artifacts, relying on evidence that includes static dependencies, code semantics, and directory organization, among other sources. Recovery produces a decomposition without a reference architecture, while mapping assigns entities to the modules of a given one. We evaluate C2A mapping, and draw on recovery research for representations of dependencies, lexical content and project organization. We organize this work by information source while keeping the two evaluation tasks distinct.

Architecture recovery techniques have historically been built around one dominant signal. Dependency-based methods treat the system as a graph and partition it to maximize modularization quality, using search-based optimization~\cite{bunch,bunch_sahc_nahc,dagc,wca,mca_eca} or clustering over dependency-derived similarity~\cite{cct,hcc_2019,limbo,fca,neighborhood_tree_2019,gma_2021}. These perform well when dependencies align with architectural boundaries, but recovery quality is sensitive to how the graph is constructed and how relations are weighted, since shared utilities and layered interactions introduce cross-module edges. Semantic methods instead cluster artifacts by identifiers and comments in a latent space learned from source code~\cite{arc,zbr,risi2012fold} or from lexical cues~\cite{kargar2017semantic}, and directory-based methods use file names and paths as an architectural signal~\cite{file_clustering}. Non-structural signals can be noisy in isolation because reused terminology and utility directories mix entities from many modules. Hybrid methods therefore combine sources through weighted similarity fusion and many-objective remodularization~\cite{sc_misra,escrc,nsga3,sadat2019multi}, consensus clustering~\cite{ensemble_recovery}, semantic-guided propagation~\cite{evol}, or weighted graphs with community detection~\cite{sade,sarif}, using different mechanisms to combine the available evidence.

Mapping shows the same pattern from a different starting point. Early work framed the assignment of unmapped entities as orphan adoption~\cite{tzerpos1997orphan}, and HuGMe introduced attraction functions to recommend mappings between entities and modules~\cite{christl2005equipping,christl2007automated}. Subsequent methods varied the signal and the learner: information-retrieval attraction based on textual similarity~\cite{bittencourt_IR}, supervised classifiers for orphan adoption~\cite{bibi2016supervised}, and a semi-supervised Naive Bayes approach combining lexical and dependency-derived information through Concrete Dependency Abstraction~\cite{olsson2019semi,olsson2022automatically}. Interactive and text-based approaches followed, comparing text classifiers for module assignment~\cite{florean2021textclassifierscomparison}, recommending mappings from natural-language architecture descriptions~\cite{InMap}, and integrating package- and class-level granularity~\cite{sinkala2023integrated}. More recent work frames method-level mapping as text classification over TF-IDF, GloVe, and LLM representations~\cite{johansson2024mapping}, examines how textual data properties condition embedding-based classification performance~\cite{johansson2025mapping}, and extracts architectural concepts from documentation and code~\cite{fuchss2025enabling}; LLM-based chain-of-thought prompting has also been used to map implementation units to a reference architecture~\cite{rukmono_deductive}, though such prompting-based approaches face scalability challenges on large codebases given context-size limits and the cost of repeatedly processing implementation artifacts.

Representation-learning approaches differ in both the information they encode and the operations used to recover or predict modules. NEGAR clusters node2vec embeddings from the dependency graph alone, with no node semantics~\cite{negar}. Jabeen et al.~\cite{structsar} build on this by adding folder organization and comparing node2vec, meta-path walks, and graph attention with and without diffusion rewiring, without lexical features. GAER and Jabeen et al.~\cite{gaer,arl_2026} instead fix folder and identifier tokens into one encoding, Word2Vec-based in the former and contrastively trained with edge attributes in the latter, without comparing it to alternatives. For mapping, prior work evaluated graph neural networks on homogeneous and dependency-type-aware graphs~\cite{jabeen2025graph}, without folder hierarchies, containment, member-level artifacts, or separating graph contribution from feature contribution. These approaches motivate our distinction between representing an information source and using it during prediction. Our configuration comparisons span lexical, structural and contextual inputs, but retain differences between their training procedures and feature content. We therefore complement them with matched dependency controls and directory-error analyses. These controls assess specific design choices; they do not turn the broader configuration comparison into an isolated ranking of information sources.
